%% ============================================================
%% sections/s360c.tex
%% United States Code, Title 21 -- § 360c. Classification of devices intended for human use
%% (FD&C Act § 513). Source identifier: /us/usc/t21/s360c
%% Relevance to the Physical AI oncology trial context:
%% Classification of devices (Class I, II, III); sets the premarket pathway and the rigor of VVUQ a bill can require of a patient-contacting oncology robot.
%% Reproduced faithfully from the OLRC USLM XML (through Pub. L. 119-93).
%% No bracketed instructions or file names are inserted into the law.
%% ============================================================
\uscsection{§~360c.}{Classification of devices intended for human use}
\uscprov{1}{\textbf{(a)}\textbf{ Classes of devices}}
\uscprov{2}{\textbf{(1)} There are established the following classes of devices intended for human use:}
\uscprov{3}{\textbf{(A)}\textbf{ Class I, General Controls.---}}
\uscprov{4}{\textbf{(i)} A device for which the controls authorized by or under section 351, 352, 360, 360f, 360h, 360i, or 360j of this title or any combination of such sections are sufficient to provide reasonable assurance of the safety and effectiveness of the device.}
\uscprov{4}{\textbf{(ii)} A device for which insufficient information exists to determine that the controls referred to in clause (i) are sufficient to provide reasonable assurance of the safety and effectiveness of the device or to establish special controls to provide such assurance, but because it---}
\uscprov{5}{\textbf{(I)} is not purported or represented to be for a use in supporting or sustaining human life or for a use which is of substantial importance in preventing impairment of human health, and}
\uscprov{5}{\textbf{(II)} does not present a potential unreasonable risk of illness or injury,}
\uscprov{4}{is to be regulated by the controls referred to in clause (i).}
\uscprov{3}{\textbf{(B)}\textbf{ Class II, Special Controls.---} A device which cannot be classified as a class I device because the general controls by themselves are insufficient to provide reasonable assurance of the safety and effectiveness of the device, and for which there is sufficient information to establish special controls to provide such assurance, including the promulgation of performance standards, postmarket surveillance, patient registries, development and dissemination of guidelines (including guidelines for the submission of clinical data in premarket notification submissions in accordance with section 360(k) of this title), recommendations, and other appropriate actions as the Secretary deems necessary to provide such assurance. For a device that is purported or represented to be for a use in supporting or sustaining human life, the Secretary shall examine and identify the special controls, if any, that are necessary to provide adequate assurance of safety and effectiveness and describe how such controls provide such assurance.}
\uscprov{3}{\textbf{(C)}\textbf{ Class III, Premarket Approval.---} A device which because---}
\uscprov{4}{\textbf{(i)} it (I) cannot be classified as a class I device because insufficient information exists to determine that the application of general controls are sufficient to provide reasonable assurance of the safety and effectiveness of the device, and (II) cannot be classified as a class II device because insufficient information exists to determine that the special controls described in subparagraph (B) would provide reasonable assurance of its safety and effectiveness, and}
\uscprov{4}{\textbf{(ii)}}
\uscprov{5}{\textbf{(I)} is purported or represented to be for a use in supporting or sustaining human life or for a use which is of substantial importance in preventing impairment of human health, or}
\uscprov{5}{\textbf{(II)} presents a potential unreasonable risk of illness or injury,}
\uscprov{3}{is to be subject, in accordance with section 360e of this title, to premarket approval to provide reasonable assurance of its safety and effectiveness.}
\uscprov{2}{If there is not sufficient information to establish a performance standard for a device to provide reasonable assurance of its safety and effectiveness, the Secretary may conduct such activities as may be necessary to develop or obtain such information.}
\uscprov{2}{\textbf{(2)} For purposes of this section and sections 360d and 360e of this title, the safety and effectiveness of a device are to be determined---}
\uscprov{3}{\textbf{(A)} with respect to the persons for whose use the device is represented or intended,}
\uscprov{3}{\textbf{(B)} with respect to the conditions of use prescribed, recommended, or suggested in the labeling of the device, and}
\uscprov{3}{\textbf{(C)} weighing any probable benefit to health from the use of the device against any probable risk of injury or illness from such use.}
\uscprov{2}{\textbf{(3)}}
\uscprov{3}{\textbf{(A)} Except as authorized by subparagraph (B), the effectiveness of a device is, for purposes of this section and sections 360d and 360e of this title, to be determined, in accordance with regulations promulgated by the Secretary, on the basis of well-controlled investigations, including 1 or more clinical investigations where appropriate, by experts qualified by training and experience to evaluate the effectiveness of the device, from which investigations it can fairly and responsibly be concluded by qualified experts that the device will have the effect it purports or is represented to have under the conditions of use prescribed, recommended, or suggested in the labeling of the device.}
\uscprov{3}{\textbf{(B)} If the Secretary determines that there exists valid scientific evidence (other than evidence derived from investigations described in subparagraph (A))---}
\uscprov{4}{\textbf{(i)} which is sufficient to determine the effectiveness of a device, and}
\uscprov{4}{\textbf{(ii)} from which it can fairly and responsibly be concluded by qualified experts that the device will have the effect it purports or is represented to have under the conditions of use prescribed, recommended, or suggested in the labeling of the device,}
\uscprov{3}{then, for purposes of this section and sections 360d and 360e of this title, the Secretary may authorize the effectiveness of the device to be determined on the basis of such evidence.}
\uscprov{3}{\textbf{(C)} In making a determination of a reasonable assurance of the effectiveness of a device for which an application under section 360e of this title has been submitted, the Secretary shall consider whether the extent of data that otherwise would be required for approval of the application with respect to effectiveness can be reduced through reliance on postmarket controls.}
\uscprov{3}{\textbf{(D)}}
\uscprov{4}{\textbf{(i)} The Secretary, upon the written request of any person intending to submit an application under section 360e of this title, shall meet with such person to determine the type of valid scientific evidence (within the meaning of subparagraphs (A) and (B)) that will be necessary to demonstrate for purposes of approval of an application the effectiveness of a device for the conditions of use proposed by such person. The written request shall include a detailed description of the device, a detailed description of the proposed conditions of use of the device, a proposed plan for determining whether there is a reasonable assurance of effectiveness, and, if available, information regarding the expected performance from the device. Within 30 days after such meeting, the Secretary shall specify in writing the type of valid scientific evidence that will provide a reasonable assurance that a device is effective under the conditions of use proposed by such person.}
\uscprov{4}{\textbf{(ii)} Any clinical data, including one or more well-controlled investigations, specified in writing by the Secretary for demonstrating a reasonable assurance of device effectiveness shall be specified as result of a determination by the Secretary that such data are necessary to establish device effectiveness. The Secretary shall consider, in consultation with the applicant, the least burdensome appropriate means of evaluating device effectiveness that would have a reasonable likelihood of resulting in approval.}
\uscprov{4}{\textbf{(iii)} For purposes of clause (ii), the term ``necessary'' means the minimum required information that would support a determination by the Secretary that an application provides reasonable assurance of the effectiveness of the device.}
\uscprov{4}{\textbf{(iv)} Nothing in this subparagraph shall alter the criteria for evaluating an application for premarket approval of a device.}
\uscprov{4}{\textbf{(v)} The determination of the Secretary with respect to the specification of valid scientific evidence under clauses (i) and (ii) shall be binding upon the Secretary, unless such determination by the Secretary could be contrary to the public health.}
\uscprov{1}{\textbf{(b)}\textbf{ Classification panels}}
\uscprov{2}{\textbf{(1)} For purposes of---}
\uscprov{3}{\textbf{(A)} determining which devices intended for human use should be subject to the requirements of general controls, performance standards, or premarket approval, and}
\uscprov{3}{\textbf{(B)} providing notice to the manufacturers and importers of such devices to enable them to prepare for the application of such requirements to devices manufactured or imported by them,}
\uscprov{2}{the Secretary shall classify all such devices (other than devices classified by subsection (f)) into the classes established by subsection (a). For the purpose of securing recommendations with respect to the classification of devices, the Secretary shall establish panels of experts or use panels of experts established before May 28, 1976, or both. Section 1013 of title 5 shall not apply to the duration of a panel established under this paragraph.}
\uscprov{2}{\textbf{(2)} The Secretary shall appoint to each panel established under paragraph (1) persons who are qualified by training and experience to evaluate the safety and effectiveness of the devices to be referred to the panel and who, to the extent feasible, possess skill in the use of, or experience in the development, manufacture, or utilization of, such devices. The Secretary shall make appointments to each panel so that each panel shall consist of members with adequately diversified expertise in such fields as clinical and administrative medicine, engineering, biological and physical sciences, and other related professions. In addition, each panel shall include as nonvoting members a representative of consumer interests and a representative of interests of the device manufacturing industry. Scientific, trade, and consumer organizations shall be afforded an opportunity to nominate individuals for appointment to the panels. No individual who is in the regular full-time employ of the United States and engaged in the administration of this chapter may be a member of any panel. The Secretary shall designate one of the members of each panel to serve as chairman thereof.}
\uscprov{2}{\textbf{(3)} Panel members (other than officers or employees of the United States), while attending meetings or conferences of a panel or otherwise engaged in its business, shall be entitled to receive compensation at rates to be fixed by the Secretary, but not at rates exceeding the daily equivalent of the rate in effect for grade GS--18 of the General Schedule, for each day so engaged, including traveltime; and while so serving away from their homes or regular places of business each member may be allowed travel expenses (including per diem in lieu of subsistence) as authorized by section 5703 of title 5, for persons in the Government service employed intermittently.}
\uscprov{2}{\textbf{(4)} The Secretary shall furnish each panel with adequate clerical and other necessary assistance.}
\uscprov{2}{\textbf{(5)}}
\uscprov{3}{\textbf{(A)} Classification panels covering each type of device shall be scheduled to meet at such times as may be appropriate for the Secretary to meet applicable statutory deadlines.}
\uscprov{3}{\textbf{(B)} When a device is specifically the subject of review by a classification panel, the Secretary shall---}
\uscprov{4}{\textbf{(i)} ensure that adequate expertise is represented on the classification panel to assess---}
\uscprov{5}{\textbf{(I)} the disease or condition which the device is intended to cure, treat, mitigate, prevent, or diagnose; and}
\uscprov{5}{\textbf{(II)} the technology of the device; and}
\uscprov{4}{\textbf{(ii)} provide an opportunity for the person whose device is specifically the subject of panel review to provide recommendations on the expertise needed among the voting members of the panel.}
\uscprov{3}{\textbf{(C)} For purposes of subparagraph (B)(i), the term ``adequate expertise'' means that the membership of the classification panel includes---}
\uscprov{4}{\textbf{(i)} two or more voting members, with a specialty or other expertise clinically relevant to the device under review; and}
\uscprov{4}{\textbf{(ii)} at least one voting member who is knowledgeable about the technology of the device.}
\uscprov{3}{\textbf{(D)} The Secretary shall provide an annual opportunity for patients, representatives of patients, and sponsors of medical devices that may be specifically the subject of a review by a classification panel to provide recommendations for individuals with appropriate expertise to fill voting member positions on classification panels.}
\uscprov{2}{\textbf{(6)}}
\uscprov{3}{\textbf{(A)} Any person whose device is specifically the subject of review by a classification panel shall have---}
\uscprov{4}{\textbf{(i)} the same access to data and information submitted to a classification panel (except for data and information that are not available for public disclosure under section 552 of title 5) as the Secretary;}
\uscprov{4}{\textbf{(ii)} the opportunity to submit, for review by a classification panel, information that is based on the data or information provided in the application submitted under section 360e of this title by the person, which information shall be submitted to the Secretary for prompt transmittal to the classification panel; and}
\uscprov{4}{\textbf{(iii)} the same opportunity as the Secretary to participate in meetings of the panel, including, subject to the discretion of the panel chairperson, by designating a representative who will be provided a time during the panel meeting to address the panel for the purpose of correcting misstatements of fact or providing clarifying information, and permitting the person or representative to call on experts within the person's organization to address such specific issues in the time provided.}
\uscprov{3}{\textbf{(B)}}
\uscprov{4}{\textbf{(i)} Any meeting of a classification panel with respect to the review of a device shall---}
\uscprov{5}{\textbf{(I)} provide adequate time for initial presentations by the person whose device is specifically the subject of such review and by the Secretary; and}
\uscprov{5}{\textbf{(II)} encourage free and open participation by all interested persons.}
\uscprov{4}{\textbf{(ii)} Following the initial presentations described in clause (i), the panel may---}
\uscprov{5}{\textbf{(I)} pose questions to a designated representative described in subparagraph (A)(iii); and}
\uscprov{5}{\textbf{(II)} consider the responses to such questions in the panel's review of the device.}
\uscprov{2}{\textbf{(7)} After receiving from a classification panel the conclusions and recommendations of the panel on a matter that the panel has reviewed, the Secretary shall review the conclusions and recommendations, shall make a final decision on the matter in accordance with section 360e(d)(2) of this title, and shall notify the affected persons of the decision in writing and, if the decision differs from the conclusions and recommendations of the panel, shall include the reasons for the difference.}
\uscprov{2}{\textbf{(8)} A classification panel under this subsection shall not be subject to the annual chartering and annual report requirements of the~\textsuperscript{1}\,{\footnotesize ~So in original. The word ``the'' probably should not appear.} chapter 10 of title 5.}
\uscprov{1}{\textbf{(c)}\textbf{ Classification panel organization and operation}}
\uscprov{2}{\textbf{(1)} The Secretary shall organize the panels according to the various fields of clinical medicine and fundamental sciences in which devices intended for human use are used. The Secretary shall refer a device to be classified under this section to an appropriate panel established or authorized to be used under subsection (b) for its review and for its recommendation respecting the classification of the device. The Secretary shall by regulation prescribe the procedure to be followed by the panels in making their reviews and recommendations. In making their reviews of devices, the panels, to the maximum extent practicable, shall provide an opportunity for interested persons to submit data and views on the classification of the devices.}
\uscprov{2}{\textbf{(2)}}
\uscprov{3}{\textbf{(A)} Upon completion of a panel's review of a device referred to it under paragraph (1), the panel shall, subject to subparagraphs (B) and (C), submit to the Secretary its recommendation for the classification of the device. Any such recommendation shall (i) contain (I) a summary of the reasons for the recommendation, (II) a summary of the data upon which the recommendation is based, and (III) an identification of the risks to health (if any) presented by the device with respect to which the recommendation is made, and (ii) to the extent practicable, include a recommendation for the assignment of a priority for the application of the requirements of section 360d or 360e of this title to a device recommended to be classified in class II or class III.}
\uscprov{3}{\textbf{(B)} A recommendation of a panel for the classification of a device in class I shall include a recommendation as to whether the device should be exempted from the requirements of section 360, 360i, or 360j(f) of this title.}
\uscprov{3}{\textbf{(C)} In the case of a device which has been referred under paragraph (1) to a panel, and which---}
\uscprov{4}{\textbf{(i)} is intended to be implanted in the human body or is purported or represented to be for a use in supporting or sustaining human life, and}
\uscprov{4}{\textbf{(ii)}}
\uscprov{5}{\textbf{(I)} has been introduced or delivered for introduction into interstate commerce for commercial distribution before May 28, 1976, or}
\uscprov{5}{\textbf{(II)} is within a type of device which was so introduced or delivered before such date and is substantially equivalent to another device within that type,}
\uscprov{3}{such panel shall recommend to the Secretary that the device be classified in class III unless the panel determines that classification of the device in such class is not necessary to provide reasonable assurance of its safety and effectiveness. If a panel does not recommend that such a device be classified in class III, it shall in its recommendation to the Secretary for the classification of the device set forth the reasons for not recommending classification of the device in such class.}
\uscprov{2}{\textbf{(3)} The panels shall submit to the Secretary within one year of the date funds are first appropriated for the implementation of this section their recommendations respecting all devices of a type introduced or delivered for introduction into interstate commerce for commercial distribution before May 28, 1976.}
\uscprov{1}{\textbf{(d)}\textbf{ Panel recommendation; publication; priorities}}
\uscprov{2}{\textbf{(1)} Upon receipt of a recommendation from a panel respecting a device, the Secretary shall publish in the Federal Register the panel's recommendation and a proposed regulation classifying such device and shall provide interested persons an opportunity to submit comments on such recommendation and the proposed regulation. After reviewing such comments, the Secretary shall, subject to paragraph (2), by regulation classify such device.}
\uscprov{2}{\textbf{(2)}}
\uscprov{3}{\textbf{(A)} A regulation under paragraph (1) classifying a device in class I shall prescribe which, if any, of the requirements of section 360, 360i, or 360j(f) of this title shall not apply to the device. A regulation which makes a requirement of section 360, 360i, or 360j(f) of this title inapplicable to a device shall be accompanied by a statement of the reasons of the Secretary for making such requirement inapplicable.}
\uscprov{3}{\textbf{(B)} A device described in subsection (c)(2)(C) shall be classified in class III unless the Secretary determines that classification of the device in such class is not necessary to provide reasonable assurance of its safety and effectiveness. A proposed regulation under paragraph (1) classifying such a device in a class other than class III shall be accompanied by a full statement of the reasons of the Secretary (and supporting documentation and data) for not classifying such device in such class and an identification of the risks to health (if any) presented by such device.}
\uscprov{2}{\textbf{(3)} In the case of devices classified in class II and devices classified under this subsection in class III and described in section 360e(b)(1) of this title the Secretary may establish priorities which, in his discretion, shall be used in applying sections 360d and 360e of this title, as appropriate, to such devices.}
\uscprov{1}{\textbf{(e)}\textbf{ Classification changes}}
\uscprov{2}{\textbf{(1)}}
\uscprov{3}{\textbf{(A)}}
\uscprov{4}{\textbf{(i)} Based on new information respecting a device, the Secretary may, upon the initiative of the Secretary or upon petition of an interested person, change the classification of such device, and revoke, on account of the change in classification, any regulation or requirement in effect under section 360d or 360e of this title with respect to such device, by administrative order published in the Federal Register following publication of a proposed reclassification order in the Federal Register, a meeting of a device classification panel described in subsection (b), and consideration of comments to a public docket, notwithstanding subchapter II of chapter 5 of title 5. The proposed reclassification order published in the Federal Register shall set forth the proposed reclassification, and a substantive summary of the valid scientific evidence concerning the proposed reclassification, including---}
\uscprov{5}{\textbf{(I)} the public health benefit of the use of the device, and the nature and, if known, incidence of the risk of the device;}
\uscprov{5}{\textbf{(II)} in the case of a reclassification from class II to class III, why general controls pursuant to subsection (a)(1)(A) and special controls pursuant to subsection (a)(1)(B) together are not sufficient to provide a reasonable assurance of safety and effectiveness for such device; and}
\uscprov{5}{\textbf{(III)} in the case of reclassification from class III to class II, why general controls pursuant to subsection (a)(1)(A) and special controls pursuant to subsection (a)(1)(B) together are sufficient to provide a reasonable assurance of safety and effectiveness for such device.}
\uscprov{4}{\textbf{(ii)} An order under this subsection changing the classification of a device from class III to class II may provide that such classification shall not take effect until the effective date of a performance standard established under section 360d of this title for such device.}
\uscprov{3}{\textbf{(B)} Authority to issue such administrative order shall not be delegated below the Director of the Center for Devices and Radiological Health, acting in consultation with the Commissioner.}
\uscprov{2}{\textbf{(2)} By an order issued under paragraph (1), the Secretary may change the classification of a device from class III---}
\uscprov{3}{\textbf{(A)} to class II if the Secretary determines that special controls would provide reasonable assurance of the safety and effectiveness of the device and that general controls would not provide reasonable assurance of the safety and effectiveness of the device, or}
\uscprov{3}{\textbf{(B)} to class I if the Secretary determines that general controls would provide reasonable assurance of the safety and effectiveness of the device.}
\uscprov{1}{\textbf{(f)}\textbf{ Initial classification and reclassification of certain devices}}
\uscprov{2}{\textbf{(1)} Any device intended for human use which was not introduced or delivered for introduction into interstate commerce for commercial distribution before May 28, 1976, is classified in class III unless---}
\uscprov{3}{\textbf{(A)} the device---}
\uscprov{4}{\textbf{(i)} is within a type of device (I) which was introduced or delivered for introduction into interstate commerce for commercial distribution before such date and which is to be classified pursuant to subsection (b), or (II) which was not so introduced or delivered before such date and has been classified in class I or II, and}
\uscprov{4}{\textbf{(ii)} is substantially equivalent to another device within such type;}
\uscprov{3}{\textbf{(B)} the Secretary in response to a petition submitted under paragraph (3) has classified such device in class I or II; or}
\uscprov{3}{\textbf{(C)} the device is classified pursuant to a request submitted under paragraph (2).}
\uscprov{2}{A device classified in class III under this paragraph shall be classified in that class until the effective date of an order of the Secretary under paragraph (2) or (3) classifying the device in class I or II.}
\uscprov{2}{\textbf{(2)}}
\uscprov{3}{\textbf{(A)}}
\uscprov{4}{\textbf{(i)} Any person who submits a report under section 360(k) of this title for a type of device that has not been previously classified under this chapter, and that is classified into class III under paragraph (1), may request, after receiving written notice of such a classification, the Secretary to classify the device.}
\uscprov{4}{\textbf{(ii)} In lieu of submitting a report under section 360(k) of this title and submitting a request for classification under clause (i) for a device, if a person determines there is no legally marketed device upon which to base a determination of substantial equivalence (as defined in subsection (i)), a person may submit a request under this clause for the Secretary to classify the device.}
\uscprov{4}{\textbf{(iii)} Upon receipt of a request under clause (i) or (ii), the Secretary shall classify the device subject to the request under the criteria set forth in subparagraphs (A) through (C) of subsection (a)(1) within 120 days.}
\uscprov{4}{\textbf{(iv)} Notwithstanding clause (iii), the Secretary may decline to undertake a classification request submitted under clause (ii) if the Secretary identifies a legally marketed device that could provide a reasonable basis for review of substantial equivalence under paragraph (1), or when the Secretary determines that the device submitted is not of low to moderate risk or that general controls would be inadequate to control the risks and special controls to mitigate the risks cannot be developed.}
\uscprov{4}{\textbf{(v)} The person submitting the request for classification under this subparagraph may recommend to the Secretary a classification for the device and shall, if recommending classification in class II, include in the request an initial draft proposal for applicable special controls, as described in subsection (a)(1)(B), that are necessary, in conjunction with general controls, to provide reasonable assurance of safety and effectiveness and a description of how the special controls provide such assurance. Any such request shall describe the device and provide detailed information and reasons for the recommended classification.}
\uscprov{3}{\textbf{(B)}}
\uscprov{4}{\textbf{(i)} The Secretary shall by written order classify the device involved. Such classification shall be the initial classification of the device for purposes of paragraph (1) and any device classified under this paragraph shall be a predicate device for determining substantial equivalence under paragraph (1).}
\uscprov{4}{\textbf{(ii)} A device that remains in class III under this subparagraph shall be deemed to be adulterated within the meaning of section 351(f)(1)(B) of this title until approved under section 360e of this title or exempted from such approval under section 360j(g) of this title.}
\uscprov{3}{\textbf{(C)} Within 30 days after the issuance of an order classifying a device under this paragraph, the Secretary shall publish a notice in the Federal Register announcing such classification.}
\uscprov{2}{\textbf{(3)}}
\uscprov{3}{\textbf{(A)} The Secretary may initiate the reclassification of a device classified into class III under paragraph (1) of this subsection or the manufacturer or importer of a device classified under paragraph (1) may petition the Secretary (in such form and manner as he shall prescribe) for the issuance of an order classifying the device in class I or class II. Within thirty days of the filing of such a petition, the Secretary shall notify the petitioner of any deficiencies in the petition which prevent the Secretary from making a decision on the petition.}
\uscprov{3}{\textbf{(B)}}
\uscprov{4}{\textbf{(i)} Upon determining that a petition does not contain any deficiency which prevents the Secretary from making a decision on the petition, the Secretary may for good cause shown refer the petition to an appropriate panel established or authorized to be used under subsection (b). A panel to which such a petition has been referred shall not later than ninety days after the referral of the petition make a recommendation to the Secretary respecting approval or denial of the petition. Any such recommendation shall contain (I) a summary of the reasons for the recommendation, (II) a summary of the data upon which the recommendation is based, and (III) an identification of the risks to health (if any) presented by the device with respect to which the petition was filed. In the case of a petition for a device which is intended to be implanted in the human body or which is purported or represented to be for a use in supporting or sustaining human life, the panel shall recommend that the petition be denied unless the panel determines that the classification in class III of the device is not necessary to provide reasonable assurance of its safety and effectiveness. If the panel recommends that such petition be approved, it shall in its recommendation to the Secretary set forth its reasons for such recommendation.}
\uscprov{4}{\textbf{(ii)} The requirements of paragraphs (1) and (2) of subsection (c) (relating to opportunities for submission of data and views and recommendations respecting priorities and exemptions from sections 360, 360i, and 360j(f) of this title) shall apply with respect to consideration by panels of petitions submitted under subparagraph (A).}
\uscprov{3}{\textbf{(C)}}
\uscprov{4}{\textbf{(i)} Within ninety days from the date the Secretary receives the recommendation of a panel respecting a petition (but not later than 210 days after the filing of such petition) the Secretary shall by order deny or approve the petition. If the Secretary approves the petition, the Secretary shall order the classification of the device into class I or class II in accordance with the criteria prescribed by subsection (a)(1)(A) or (a)(1)(B). In the case of a petition for a device which is intended to be implanted in the human body or which is purported or represented to be for a use in supporting or sustaining human life, the Secretary shall deny the petition unless the Secretary determines that the classification in class III of the device is not necessary to provide reasonable assurance of its safety and effectiveness. An order approving such petition shall be accompanied by a full statement of the reasons of the Secretary (and supporting documentation and data) for approving the petition and an identification of the risks to health (if any) presented by the device to which such order applies.}
\uscprov{4}{\textbf{(ii)} The requirements of paragraphs (1) and (2)(A) of subsection (d) (relating to publication of recommendations, opportunity for submission of comments, and exemption from sections 360, 360i, and 360j(f) of this title) shall apply with respect to action by the Secretary on petitions submitted under subparagraph (A).}
\uscprov{2}{\textbf{(4)} If a manufacturer reports to the Secretary under section 360(k) of this title that a device is substantially equivalent to another device---}
\uscprov{3}{\textbf{(A)} which the Secretary has classified as a class III device under subsection (b),}
\uscprov{3}{\textbf{(B)} which was introduced or delivered for introduction into interstate commerce for commercial distribution before December 1, 1990, and}
\uscprov{3}{\textbf{(C)} for which no final regulation requiring premarket approval has been promulgated under section 360e(b) of this title,}
\uscprov{2}{the manufacturer shall certify to the Secretary that the manufacturer has conducted a reasonable search of all information known or otherwise available to the manufacturer respecting such other device and has included in the report under section 360(k) of this title a summary of and a citation to all adverse safety and effectiveness data respecting such other device and respecting the device for which the section 360(k) report is being made and which has not been submitted to the Secretary under section 360i of this title. The Secretary may require the manufacturer to submit the adverse safety and effectiveness data described in the report.}
\uscprov{2}{\textbf{(5)} The Secretary may not withhold a determination of the initial classification of a device under paragraph (1) because of a failure to comply with any provision of this chapter unrelated to a substantial equivalence decision, including a finding that the facility in which the device is manufactured is not in compliance with good manufacturing requirements as set forth in regulations of the Secretary under section 360j(f) of this title (other than a finding that there is a substantial likelihood that the failure to comply with such regulations will potentially present a serious risk to human health).}
\uscprov{2}{\textbf{(6)}}
\uscprov{3}{\textbf{(A)} Subject to the succeeding subparagraphs of this paragraph, the Secretary shall, by written order, classify an accessory under this section based on the risks of the accessory when used as intended and the level of regulatory controls necessary to provide a reasonable assurance of safety and effectiveness of the accessory, notwithstanding the classification of any other device with which such accessory is intended to be used.}
\uscprov{3}{\textbf{(B)} The classification of any accessory distinct from another device by regulation or written order issued prior to December 13, 2016, shall continue to apply unless and until the accessory is reclassified by the Secretary, notwithstanding the classification of any other device with which such accessory is intended to be used. Nothing in this paragraph shall preclude the Secretary's authority to initiate the classification of an accessory through regulation or written order, as appropriate.}
\uscprov{3}{\textbf{(C)}}
\uscprov{4}{\textbf{(i)} In the case of a device intended to be used with an accessory, where the accessory has been included in an application for premarket approval of such device under section 360e of this title or a report under section 360(k) of this title for clearance of such device and the Secretary has not classified such accessory distinctly from another device in accordance with subparagraph (A), the person filing the application or report (as applicable) at the time such application or report is filed---}
\uscprov{5}{\textbf{(I)} may include a written request for the proper classification of the accessory pursuant to subparagraph (A);}
\uscprov{5}{\textbf{(II)} shall include in any such request such information as may be necessary for the Secretary to evaluate, based on the least burdensome approach, the appropriate class for the accessory under subsection (a); and}
\uscprov{5}{\textbf{(III)} shall, if the request under subclause (I) is requesting classification of the accessory in class II, include in the application an initial draft proposal for special controls, if special controls would be required pursuant to subsection (a)(1)(B).}
\uscprov{4}{\textbf{(ii)} The Secretary's response under section 360e(d) or section 360(n) of this title (as applicable) to an application or report described in clause (i) shall also contain the Secretary's granting or denial of the request for classification of the accessory involved.}
\uscprov{4}{\textbf{(iii)} The Secretary's evaluation of an accessory under clause (i) shall constitute an order establishing a new classification for such accessory for the specified intended use or uses of such accessory and for any accessory with the same intended use or uses as such accessory.}
\uscprov{3}{\textbf{(D)} For accessories that have been granted marketing authorization as part of a submission for another device with which the accessory involved is intended to be used, through an application for such other device under section 360e(c) of this title, a report under section 360(k) of this title, or a request for classification under paragraph (2) of this subsection, the following shall apply:}
\uscprov{4}{\textbf{(i)} Not later than the date that is one year after August 18, 2017, and at least once every 5 years thereafter, and as the Secretary otherwise determines appropriate, pursuant to this paragraph, the Secretary shall publish in the Federal Register a notice proposing a list of such accessories that the Secretary determines may be suitable for a distinct classification in class I and the proposed regulations for such classifications. In developing such list, the Secretary shall consider recommendations from sponsors of device submissions and other stakeholders for accessories to be included on such list. The notices shall provide for a period of not less than 60 calendar days for public comment. Within 180 days after the end of the comment period, the Secretary shall publish in the Federal Register a final action classifying such suitable accessories into class I.}
\uscprov{4}{\textbf{(ii)} A manufacturer or importer of an accessory that has been granted such marketing authorization may submit to the Secretary a written request for the appropriate classification of the accessory based on the risks and appropriate level of regulatory controls as described in subparagraph (A), and shall, if the request is requesting classification of the accessory in class II, include in the submission an initial draft proposal for special controls, if special controls would be required pursuant to subsection (a)(1)(B). Such request shall include such information as may be necessary for the Secretary to evaluate, based on the least burdensome approach, the appropriate class for the accessory under subsection (a). The Secretary shall provide an opportunity for a manufacturer or importer to meet with appropriate personnel of the Food and Drug Administration to discuss the appropriate classification of such accessory prior to submitting a written request under this clause for classification of the accessory.}
\uscprov{4}{\textbf{(iii)} The Secretary shall respond to a request made under clause (ii) not later than 85 calendar days after receiving such request by issuing a written order classifying the accessory or denying the request. If the Secretary does not agree with the recommendation for classification submitted by the manufacturer or importer, the response shall include a detailed description and justification for such determination. Within 30 calendar days after granting such a request, the Secretary shall publish a notice in the Federal Register announcing such response.}
\uscprov{3}{\textbf{(E)} Nothing in this paragraph may be construed as precluding a manufacturer of an accessory of a new type from using the classification process described in subsection (f)(2) to obtain classification of such accessory in accordance with the criteria and requirements set forth in that subsection.}
\uscprov{1}{\textbf{(g)}\textbf{ Information} Within sixty days of the receipt of a written request of any person for information respecting the class in which a device has been classified or the requirements applicable to a device under this chapter, the Secretary shall provide such person a written statement of the classification (if any) of such device and the requirements of this chapter applicable to the device.}
\uscprov{1}{\textbf{(h)}\textbf{ Definitions} For purposes of this section and sections 351, 360, 360d, 360e, 360f, 360i, and 360j of this title}
\uscprov{2}{\textbf{(1)} a reference to ``general controls'' is a reference to the controls authorized by or under sections 351, 352, 360, 360f, 360h, 360i, and 360j of this title,}
\uscprov{2}{\textbf{(2)} a reference to ``class I'', ``class II'', or ``class III'' is a reference to a class of medical devices described in subparagraph (A), (B), or (C) of subsection (a)(1), and}
\uscprov{2}{\textbf{(3)} a reference to a ``panel under section 360c of this title'' is a reference to a panel established or authorized to be used under this section.}
\uscprov{1}{\textbf{(i)}\textbf{ Substantial equivalence}}
\uscprov{2}{\textbf{(1)}}
\uscprov{3}{\textbf{(A)} For purposes of determinations of substantial equivalence under subsection (f) and section 360j(\textit{l}) of this title, the term ``substantially equivalent'' or ``substantial equivalence'' means, with respect to a device being compared to a predicate device, that the device has the same intended use as the predicate device and that the Secretary by order has found that the device---}
\uscprov{4}{\textbf{(i)} has the same technological characteristics as the predicate device, or}
\uscprov{4}{\textbf{(ii)}}
\uscprov{5}{\textbf{(I)} has different technological characteristics and the information submitted that the device is substantially equivalent to the predicate device contains information, including appropriate clinical or scientific data if deemed necessary by the Secretary or a person accredited under section 360m of this title, that demonstrates that the device is as safe and effective as a legally marketed device, and (II) does not raise different questions of safety and effectiveness than the predicate device.}
\uscprov{3}{\textbf{(B)} For purposes of subparagraph (A), the term ``different technological characteristics'' means, with respect to a device being compared to a predicate device, that there is a significant change in the materials, design, energy source, or other features of the device from those of the predicate device.}
\uscprov{3}{\textbf{(C)} To facilitate reviews of reports submitted to the Secretary under section 360(k) of this title, the Secretary shall consider the extent to which reliance on postmarket controls may expedite the classification of devices under subsection (f)(1) of this section.}
\uscprov{3}{\textbf{(D)}}
\uscprov{4}{\textbf{(i)} Whenever the Secretary requests information to demonstrate that devices with differing technological characteristics are substantially equivalent, the Secretary shall only request information that is necessary to making substantial equivalence determinations. In making such request, the Secretary shall consider the least burdensome means of demonstrating substantial equivalence and request information accordingly.}
\uscprov{4}{\textbf{(ii)} For purposes of clause (i), the term ``necessary'' means the minimum required information that would support a determination of substantial equivalence between a new device and a predicate device.}
\uscprov{4}{\textbf{(iii)} Nothing in this subparagraph shall alter the standard for determining substantial equivalence between a new device and a predicate device.}
\uscprov{3}{\textbf{(E)}}
\uscprov{4}{\textbf{(i)} Any determination by the Secretary of the intended use of a device shall be based upon the proposed labeling submitted in a report for the device under section 360(k) of this title. However, when determining that a device can be found substantially equivalent to a legally marketed device, the director of the organizational unit responsible for regulating devices (in this subparagraph referred to as the ``Director'') may require a statement in labeling that provides appropriate information regarding a use of the device not identified in the proposed labeling if, after providing an opportunity for consultation with the person who submitted such report, the Director determines and states in writing---}
\uscprov{5}{\textbf{(I)} that there is a reasonable likelihood that the device will be used for an intended use not identified in the proposed labeling for the device; and}
\uscprov{5}{\textbf{(II)} that such use could cause harm.}
\uscprov{4}{\textbf{(ii)} Such determination shall---}
\uscprov{5}{\textbf{(I)} be provided to the person who submitted the report within 10 days from the date of the notification of the Director's concerns regarding the proposed labeling;}
\uscprov{5}{\textbf{(II)} specify the limitations on the use of the device not included in the proposed labeling; and}
\uscprov{5}{\textbf{(III)} find the device substantially equivalent if the requirements of subparagraph (A) are met and if the labeling for such device conforms to the limitations specified in subclause (II).}
\uscprov{4}{\textbf{(iii)} The responsibilities of the Director under this subparagraph may not be delegated.}
\uscprov{3}{\textbf{(F)} Not later than 270 days after November 21, 1997, the Secretary shall issue guidance specifying the general principles that the Secretary will consider in determining when a specific intended use of a device is not reasonably included within a general use of such device for purposes of a determination of substantial equivalence under subsection (f) or section 360j(\textit{l}) of this title.}
\uscprov{2}{\textbf{(2)} A device may not be found to be substantially equivalent to a predicate device that has been removed from the market at the initiative of the Secretary or that has been determined to be misbranded or adulterated by a judicial order.}
\uscprov{2}{\textbf{(3)}}
\uscprov{3}{\textbf{(A)} As part of a submission under section 360(k) of this title respecting a device, the person required to file a premarket notification under such section shall provide an adequate summary of any information respecting safety and effectiveness or state that such information will be made available upon request by any person.}
\uscprov{3}{\textbf{(B)} Any summary under subparagraph (A) respecting a device shall contain detailed information regarding data concerning adverse health effects and shall be made available to the public by the Secretary within 30 days of the issuance of a determination that such device is substantially equivalent to another device.}
\uscprov{1}{\textbf{(j)}\textbf{ Training and oversight of least burdensome requirements}}
\uscprov{2}{\textbf{(1)} The Secretary shall---}
\uscprov{3}{\textbf{(A)} ensure that each employee of the Food and Drug Administration who is involved in the review of premarket submissions, including supervisors, receives training regarding the meaning and implementation of the least burdensome requirements under subsections (a)(3)(D) and (i)(1)(D) of this section and section 360e(c)(5) of this title; and}
\uscprov{3}{\textbf{(B)} periodically assess the implementation of the least burdensome requirements, including the employee training under subparagraph (A), to ensure that the least burdensome requirements are fully and consistently applied.}
\uscprov{2}{\textbf{(2)} Not later than 18 months after December 13, 2016, the ombudsman for any organizational unit of the Food and Drug Administration responsible for the premarket review of devices shall---}
\uscprov{3}{\textbf{(A)} conduct an audit of the training described in paragraph (1)(A), including the effectiveness of such training in implementing the least burdensome requirements;}
\uscprov{3}{\textbf{(B)} include in such audit interviews of persons who are representatives of the device industry regarding their experiences in the device premarket review process, including with respect to the application of least burdensome concepts to premarket review and decisionmaking;}
\uscprov{3}{\textbf{(C)} include in such audit a list of the measurement tools the Secretary uses to assess the implementation of the least burdensome requirements, including under paragraph (1)(B) and section 360g--1(a)(3) of this title, and may also provide feedback on the effectiveness of such tools in the implementation of the least burdensome requirements;}
\uscprov{3}{\textbf{(D)} summarize the findings of such audit in a final audit report; and}
\uscprov{3}{\textbf{(E)} within 30 calendar days of completion of such final audit report, make such final audit report available---}
\uscprov{4}{\textbf{(i)} to the Committee on Health, Education, Labor, and Pensions of the Senate and the Committee on Energy and Commerce of the House of Representatives; and}
\uscprov{4}{\textbf{(ii)} on the Internet website of the Food and Drug Administration.}
\uscprov{1}{\textbf{(k)}\textbf{ Dual submission for certain devices} For a device authorized for emergency use under section 360bbb--3 of this title for which, in accordance with section 360bbb--3(m) of this title, the Secretary has deemed a laboratory examination or procedure associated with such device to be in the category of examinations and procedures described in section 263a(d)(3) of title 42, the sponsor of such device may, when submitting a request for classification under subsection (f)(2), submit a single submission containing---}
\uscprov{2}{\textbf{(1)} the information needed for such a request; and}
\uscprov{2}{\textbf{(2)} sufficient information to enable the Secretary to determine whether such laboratory examination or procedure satisfies the criteria to be categorized under section 263a(d)(3) of title 42.}
\uscsource{(June 25, 1938, ch. 675, §~513, as added Pub. L. 94--295, §~2, May 28, 1976, 90 Stat. 540; amended Pub. L. 101--629, §§~4(a), 5(a)--(c)(1), (3), 12(a), 18(a), Nov. 28, 1990, 104 Stat. 4515, 4517, 4518, 4523, 4528; Pub. L. 102--300, §~6(e), June 16, 1992, 106 Stat. 240; Pub. L. 103--80, §~3(s), Aug. 13, 1993, 107 Stat. 778; Pub. L. 105--115, title II, §§~205(a), (b), 206(b), (c), 207, 208, 217, Nov. 21, 1997, 111 Stat. 2336, 2337, 2339, 2340, 2350; Pub. L. 107--250, title II, §~208, Oct. 26, 2002, 116 Stat. 1613; Pub. L. 112--144, title VI, §§~602, 607--608(a)(2)(A), July 9, 2012, 126 Stat. 1051, 1054--1056; Pub. L. 114--255, div. A, title III, §§~3055, 3058(a), 3060(c), 3101(a)(2)(I), Dec. 13, 2016, 130 Stat. 1127, 1128, 1133, 1154; Pub. L. 115--52, title VII, §~707(a), (b), title IX, §~901(h), Aug. 18, 2017, 131 Stat. 1060, 1062, 1077; Pub. L. 117--286, §~4(a)(155), Dec. 27, 2022, 136 Stat. 4323; Pub. L. 117--328, div. FF, title III, §~3301, Dec. 29, 2022, 136 Stat. 5831.)}
\usccrosshead{Editorial Notes}
\uscnotehead{Amendments}
\uscnotepar{2022---Subsec. (b)(1). Pub. L. 117--286, §~4(a)(155)(A), substituted ``Section 1013 of title 5'' for ``Section 14 of the Federal Advisory Committee Act'' in concluding provisions.}
\uscnotepar{Subsec. (b)(8). Pub. L. 117--286, §~4(a)(155)(B), substituted ``chapter 10 of title 5.'' for ``Federal Advisory Committee Act.''}
\uscnotepar{Subsec. (k). Pub. L. 117--328 added subsec. (k).}
\uscnotepar{2017---Subsec. (b)(5)(D). Pub. L. 115--52, §~901(h), substituted ``medical devices that may be specifically the subject of a review by a classification panel'' for ``medical device submissions''.}
\uscnotepar{Subsec. (b)(9). Pub. L. 115--52, §~707(b), struck out par. (9) which read as follows: ``The Secretary shall classify an accessory under this section based on the intended use of the accessory, notwithstanding the classification of any other device with which such accessory is intended to be used.''}
\uscnotepar{Subsec. (f)(6). Pub. L. 115--52, §~707(a), added par. (6).}
\uscnotepar{2016---Subsec. (b)(5). Pub. L. 114--255, §~3055(a), designated existing provisions as subpar. (A) and added subpars. (B) to (D).}
\uscnotepar{Subsec. (b)(6)(A)(iii). Pub. L. 114--255, §~3055(b)(1), inserted before period at end ``, including, subject to the discretion of the panel chairperson, by designating a representative who will be provided a time during the panel meeting to address the panel for the purpose of correcting misstatements of fact or providing clarifying information, and permitting the person or representative to call on experts within the person's organization to address such specific issues in the time provided''.}
\uscnotepar{Subsec. (b)(6)(B). Pub. L. 114--255, §~3055(b)(2), added subpar. (B) and struck out former subpar. (B) which read as follows: ``Any meetings of a classification panel shall provide adequate time for initial presentations and for response to any differing views by persons whose devices are specifically the subject of a classification panel review, and shall encourage free and open participation by all interested persons.''}
\uscnotepar{Subsec. (b)(9). Pub. L. 114--255, §~3060(c), added par. (9).}
\uscnotepar{Subsec. (f)(2)(A)(i). Pub. L. 114--255, §~3101(a)(2)(I)(i), struck out ``within 30 days'' after ``may request,''.}
\uscnotepar{Subsec. (f)(2)(A)(iv). Pub. L. 114--255, §~3101(a)(2)(I)(ii), substituted ``low to moderate'' for ``low-moderate''.}
\uscnotepar{Subsec. (j). Pub. L. 114--255, §~3058(a), added subsec. (j).}
\uscnotepar{2012---Subsec. (a)(3)(D)(iii) to (v). Pub. L. 112--144, §~602(a), added cls. (iii) and (iv) and redesignated former cl. (iii) as (v).}
\uscnotepar{Subsec. (e)(1). Pub. L. 112--144, §~608(a)(1), amended par. (1) generally. Prior to amendment, par. (1) read as follows: ``Based on new information respecting a device, the Secretary may, upon his own initiative or upon petition of an interested person, by regulation (A) change such device's classification, and (B) revoke, because of the change in classification, any regulation or requirement in effect under section 360d or 360e of this title with respect to such device. In the promulgation of such a regulation respecting a device's classification, the Secretary may secure from the panel to which the device was last referred pursuant to subsection (c) of this section a recommendation respecting the proposed change in the device's classification and shall publish in the Federal Register any recommendation submitted to the Secretary by the panel respecting such change. A regulation under this subsection changing the classification of a device from class III to class II may provide that such classification shall not take effect until the effective date of a performance standard established under section 360d of this title for such device.''}
\uscnotepar{Subsec. (e)(2). Pub. L. 112--144, §~608(a)(2)(A), substituted ``an order issued'' for ``regulation promulgated'' in introductory provisions.}
\uscnotepar{Subsec. (f)(1)(C). Pub. L. 112--144, §~607(b), added subpar. (C).}
\uscnotepar{Subsec. (f)(2)(A). Pub. L. 112--144, §~607(a)(1)--(3), designated existing provisions as cl. (i), struck out ``under the criteria set forth in subparagraphs (A) through (C) of subsection (a)(1) of this section. The person may, in the request, recommend to the Secretary a classification for the device. Any such request shall describe the device and provide detailed information and reasons for the recommended classification'' before period at end, and added cls. (ii) to (v).}
\uscnotepar{Subsec. (f)(2)(B)(i). Pub. L. 112--144, §~607(a)(4), substituted ``The Secretary'' for ``Not later than 60 days after the date of the submission of the request under subparagraph (A), the Secretary''.}
\uscnotepar{Subsec. (i)(1)(D). Pub. L. 112--144, §~602(b), designated existing provisions as cl. (i) and added cls. (ii) and (iii).}
\uscnotepar{2002---Subsec. (i)(1)(E)(iv). Pub. L. 107--250 struck out cl. (iv) which read as follows: ``This subparagraph has no legal effect after the expiration of the five-year period beginning on November 21, 1997.''}
\uscnotepar{1997---Subsec. (a)(3)(A). Pub. L. 105--115, §~217, substituted ``1 or more clinical investigations'' for ``clinical investigations''.}
\uscnotepar{Subsec. (a)(3)(C), (D). Pub. L. 105--115, §~205(a), added subpars. (C) and (D).}
\uscnotepar{Subsec. (b)(5) to (8). Pub. L. 105--115, §~208, added pars. (5) to (8).}
\uscnotepar{Subsec. (f)(1). Pub. L. 105--115, §~207(1)(B), substituted ``paragraph (2) or (3)'' for ``paragraph (2)'' in closing provisions.}
\uscnotepar{Subsec. (f)(1)(B). Pub. L. 105--115, §~207(1)(A), substituted ``paragraph (3)'' for ``paragraph (2)''.}
\uscnotepar{Subsec. (f)(2) to (4). Pub. L. 105--115, §~207(2), (3), added par. (2) and redesignated former pars. (2) and (3) as (3) and (4), respectively.}
\uscnotepar{Subsec. (f)(5). Pub. L. 105--115, §~206(b), added par. (5).}
\uscnotepar{Subsec. (i)(1)(A)(ii). Pub. L. 105--115, §~206(c)(1), substituted ``appropriate clinical or scientific data'' for ``clinical data'', inserted ``or a person accredited under section 360m of this title'' after ``Secretary'', and substituted ``effectiveness'' for ``efficacy''.}
\uscnotepar{Subsec. (i)(1)(C) to (E). Pub. L. 105--115, §~205(b), added subpars. (C) to (E).}
\uscnotepar{Subsec. (i)(1)(F). Pub. L. 105--115, §~206(c)(2), added subpar. (F).}
\uscnotepar{1993---Subsec. (b)(3). Pub. L. 103--80 substituted ``5703'' for ``5703(b)''.}
\uscnotepar{1992---Subsec. (f)(3). Pub. L. 102--300 redesignated clauses (i) to (iii) as subpars. (A) to (C), respectively, and substituted ``the section 360(k) report'' for ``the 360(k) report'' in closing provisions.}
\uscnotepar{1990---Subsec. (a)(1)(A)(ii). Pub. L. 101--629, §~5(a)(1), substituted ``or to establish special controls'' for ``or to establish a performance standard''.}
\uscnotepar{Subsec. (a)(1)(B). Pub. L. 101--629, §~5(a)(2), amended subpar. (B) generally. Prior to amendment, subpar. (B) read as follows: ``Class II, Performance Standards.---A device which cannot be classified as a class I device because the controls authorized by or under sections 351, 352, 360, 360f, 360h, 360i, and 360j of this title by themselves are insufficient to provide reasonable assurance of the safety and effectiveness of the device, for which there is sufficient information to establish a performance standard to provide such assurance, and for which it is therefore necessary to establish for the device a performance standard under section 360d of this title to provide reasonable assurance of its safety and effectiveness.''}
\uscnotepar{Subsec. (a)(1)(C)(i). Pub. L. 101--629, §~5(a)(3), amended cl. (i) generally. Prior to amendment, cl. (i) read as follows: ``it (I) cannot be classified as a class I device because insufficient information exists to determine that the controls authorized by or under sections 351, 352, 360, 360f, 360h, 360i, and 360j of this title are sufficient to provide reasonable assurance of the safety and effectiveness of the device and (II) cannot be classified as a class II device because insufficient information exists for the establishment of a performance standard to provide reasonable assurance of its safety and effectiveness, and''.}
\uscnotepar{Subsec. (e). Pub. L. 101--629, §~5(b), designated existing provisions as par. (1), redesignated cls. (1) and (2) as (A) and (B), respectively, and added par. (2).}
\uscnotepar{Subsec. (f). Pub. L. 101--629, §~5(c)(3), inserted ``and reclassification'' before ``of'' in heading.}
\uscnotepar{Subsec. (f)(2)(A). Pub. L. 101--629, §~5(c)(1), substituted ``The Secretary may initiate the reclassification of a device classified into class III under paragraph (1) of this subsection or the manufacturer'' for ``The manufacturer''.}
\uscnotepar{Subsec. (f)(2)(B)(i). Pub. L. 101--629, §~18(a), substituted ``the Secretary may for good cause shown'' for ``the Secretary shall''.}
\uscnotepar{Subsec. (f)(3). Pub. L. 101--629, §~4(a), added par. (3).}
\uscnotepar{Subsec. (i). Pub. L. 101--629, §~12(a), added subsec. (i).}
\usccrosshead{Statutory Notes and Related Subsidiaries}
\uscnotehead{Effective Date of 2017 Amendment}
\uscnotepar{Pub. L. 115--52, title VII, §~707(c), Aug. 18, 2017, 131 Stat. 1062, provided that: ``The amendments made by subsections (a) and (b) [amending this section] shall take effect on the date that is 60 days after the date of enactment of this Act [Aug. 18, 2017].''}
\uscnotehead{Effective Date of 1997 Amendment}
\uscnotepar{Amendment by Pub. L. 105--115 effective 90 days after Nov. 21, 1997, except as otherwise provided, see section 501 of Pub. L. 105--115, set out as a note under section 321 of this title.}
\uscnotehead{Short Title of 1976 Amendment}
\uscnotepar{Pub. L. 94--295, §~1(a), May 28, 1976, 90 Stat. 539, provided that: ``This Act [enacting sections 360c to 360k, 379, and 379a of this title and section 3512 of Title 42, The Public Health and Welfare, and amending sections 321, 331, 334, 351, 352, 358, 360, 374, 379e, and 381 of this title and section 55 of Title 15, Commerce and Trade] may be cited as the `Medical Device Amendments of 1976'.''}
\uscnotehead{Regulations}
\uscnotepar{Pub. L. 101--629, §~12(b), Nov. 28, 1990, 104 Stat. 4524, provided that: ``Within 12 months of the date of the enactment of this Act [Nov. 28, 1990], the Secretary of Health and Human Services shall issue regulations establishing the requirements of the summaries under section 513(i)(3) of the Federal Food, Drug, and Cosmetic Act [21 U.S.C. 360c(i)(3)], as added by the amendment made by subsection (a).''}
\uscnotehead{Devices Reclassified Prior to July 9, 2012}
\uscnotepar{Pub. L. 112--144, title VI, §~608(a)(3), July 9, 2012, 126 Stat. 1056, provided that: ``(A) In general.---The amendments made by this subsection [amending this section and sections 360d and 360g of this title] shall have no effect on a regulation promulgated with respect to the classification of a device under section 513(e) of the Federal Food, Drug, and Cosmetic Act [21 U.S.C. 360c(e)] prior to the date of enactment of this Act [July 9, 2012]. ``(B) Applicability of other provisions.---In the case of a device reclassified under section 513(e) of the Federal Food, Drug, and Cosmetic Act [21 U.S.C. 360c(e)] by regulation prior to the date of enactment of this Act [July 9, 2012], section 517(a)(1) of the Federal Food, Drug, and Cosmetic Act (21 U.S.C. 360g(a)(1)) shall apply to such regulation promulgated under section 513(e) of such Act with respect to such device in the same manner such section 517(a)(1) applies to an administrative order issued with respect to a device reclassified after the date of enactment of this Act.''}
\uscnotehead{Daily Wear Soft or Daily Wear Nonhydrophilic Plastic Contact Lenses}
\uscnotepar{Pub. L. 101--629, §~4(b)(3), Nov. 28, 1990, 104 Stat. 4517, provided that: ``(A) Notwithstanding section 520(\textit{l})(5) of the Federal Food, Drug, and Cosmetic Act [21 U.S.C. 360j(\textit{l})(5)], the Secretary of Health and Human Services shall not retain any daily wear soft or daily wear nonhydrophilic plastic contact lens in class III under such Act [this chapter] unless the Secretary finds that it meets the criteria set forth in section 513(a)(1)(C) of such Act [21 U.S.C. 360c(a)(1)(C)]. The finding and the grounds for the finding shall be published in the Federal Register. For any such lens, the Secretary shall make the determination respecting reclassification required in section 520(\textit{l})(5)(B) of such Act within 24 months of the date of the enactment of this paragraph [Nov. 28, 1990]. ``(B) The Secretary of Health and Human Services may by notice published in the Federal Register extend the two-year period prescribed by subparagraph (A) for a lens for an additional period not to exceed one year. ``(C)(i) Before classifying a lens in class II pursuant to subparagraph (A), the Secretary of Health and Human Services shall pursuant to section 513(a)(1)(B) of such Act assure that appropriate regulatory safeguards are in effect which provide reasonable assurance of the safety and effectiveness of such lens, including clinical and preclinical data if deemed necessary by the Secretary. ``(ii) Prior to classifying a lens in class I pursuant to subparagraph (A), the Secretary shall assure that appropriate regulatory safeguards are in effect which provide reasonable assurance of the safety and effectiveness of such lens, including clinical and preclinical data if deemed necessary by the Secretary. ``(D) Notwithstanding section 520(\textit{l})(5) of such Act, if the Secretary of Health and Human Services has not made the finding and published the finding required by subparagraph (A) within 36 months of the date of the enactment of this subparagraph [Nov. 28, 1990], the Secretary shall issue an order placing the lens in class II. ``(E) Any person adversely affected by a final regulation under this paragraph revising the classification of a lens may challenge the revision of the classification of such lens only by filing a petition under section 513(e) for a classification change.''}
\uscnotehead{References in Other Laws to GS--16, 17, or 18 Pay Rates}
\uscnotepar{References in laws to the rates of pay for GS--16, 17, or 18, or to maximum rates of pay under the General Schedule, to be considered references to rates payable under specified sections of Title 5, Government Organization and Employees, see section 529 [title I, §~101(c)(1)] of Pub. L. 101--509, set out in a note under section 5376 of Title 5.}
